Demanda:
Cantidad de productos requeridos durante un período determinado, calculada a partir del comportamiento histórico de consumo o ventas.

Predicción de demanda:
Proceso mediante el cual se estiman las cantidades futuras de productos utilizando modelos estadísticos o de aprendizaje automático basados en datos históricos.

Gestión de inventarios:
Conjunto de actividades orientadas al control, almacenamiento, abastecimiento y seguimiento de los productos disponibles en una organización.

Inventario
Conjunto de productos, materiales o recursos almacenados y disponibles para su uso, venta o distribución.

Desabastecimiento:
Situación en la que un producto no se encuentra disponible en el inventario cuando es requerido.

Sistema de apoyo a la decisión:
Herramienta informática diseñada para asistir a los usuarios en la toma de decisiones mediante el análisis de datos y la generación de información relevante.

Modelo de predicción:
Algoritmo o método matemático utilizado para estimar valores futuros a partir de datos históricos y patrones identificados en la información.

Datos históricos:
Registros de consumo, ventas, compras o movimientos de inventario ocurridos en períodos anteriores y utilizados para el análisis y la predicción.

Punto de reposición:
Nivel de inventario a partir del cual se debe realizar un nuevo pedido para evitar el desabastecimiento.

Tiempo de entrega (Lead Time):
Período que transcurre entre la realización de un pedido y la recepción efectiva de los productos solicitados.

Series de tiempo:
Conjunto de observaciones registradas en orden cronológico que permiten analizar patrones, tendencias y comportamientos a lo largo del tiempo.

Aprendizaje automático (Machine Learning):
Rama de la inteligencia artificial que permite a los sistemas aprender patrones a partir de datos y realizar predicciones sin ser programados explícitamente para cada situación.

Dominio de inventario:
Contexto o sector específico en el que se gestionan inventarios, como retail, manufactura, logística, salud, restaurantes o comercio electrónico.